
\documentclass{beamer}
\usepackage[french]{babel}
\usepackage[utf8]{inputenc}
\usepackage[T1]{fontenc}

% Thème de la présentation
\usetheme{Madrid}
\usecolortheme{whale}

\title[Génération Automatique de Code GUI]{Modèle de Génération Automatique de Code à partir de Maquettes GUI Haute Fidélité}
\subtitle{Utilisation des techniques de Deep Learning}
\author{Résumé de l'article de Samir et al. (2024)}
\date{\today}

\begin{document}

% Slide de titre
\begin{frame}
    \titlepage
\end{frame}

% --- INTRODUCTION ---
\begin{frame}{Introduction}
    \begin{itemize}
        \item Les interfaces graphiques (GUI) sont essentielles pour l'interaction utilisateur-logiciel en raison de leur nature visuelle .
        \item Le flux de travail classique repose sur la création de \textbf{maquettes (mockups)} de haute fidélité avant l'implémentation.
        \item Le développement des GUI est souvent une tâche \textbf{manuelle et chronophage} .
        \item Environ 45\% de la taille du code d'un logiciel concerne l'interface utilisateur, occupant près de 50\% du temps de développement.
    \end{itemize}
\end{frame}

% --- LE PROBLEME ---
\begin{frame}{Le Problème}
    \begin{block}{Constats majeurs}
        \begin{itemize}
            \item Processus fastidieux et répétitif dû aux changements fréquents de spécifications .
            \item Les développeurs passent trop de temps sur l'UI au détriment de la logique métier du logiciel.
            \item Tâche sujette aux erreurs lors de l'extraction des éléments visuels et de leur hiérarchie .
            \item Nécessité de réimplémenter l'interface pour chaque plateforme native (Android, iOS, Web) .
        \end{itemize}
    \end{block}
\end{frame}

% --- PROBLEMATIQUE ---
\begin{frame}{Problématique}
    \begin{center}
        \Large \textbf{Comment automatiser efficacement la conversion d'une image de maquette haute fidélité en code frontal multiplateforme tout en conservant la hiérarchie visuelle ?}
    \end{center}
    \vspace{0.5cm}
    \begin{itemize}
        \item Compréhension visuelle par la machine (Computer Vision) .
        \item Diversité des designs et complexité des composants GUI.
        \item Manque de précision des jeux de données existants (Rico, etc.).
    \end{itemize}
\end{frame}

% --- SOLUTIONS APPORTÉES (Architecture) ---
\begin{frame}{Solutions apportées : Architecture}
    L'étude propose un pipeline en cinq phases  :
    \begin{enumerate}
        \item \textbf{Préparation du modèle :} Utilisation de YOLOv7 avec réglage d'hyperparamètres (K-Mean++ pour les anchor boxes).
        \item \textbf{Détection d'objets :} Séparation en deux modèles distincts pour améliorer la précision :
            \begin{itemize}
                \item \textbf{Modèle Atomique :} (boutons, icônes, etc.) .
                \item \textbf{Modèle Conteneur :} (barres d'outils, tiroirs, etc.) .
            \end{itemize}
        \item \textbf{Post-traitement :} Extraction des styles (couleurs, tailles) et du texte via OCR (Tesseract).
        \item \textbf{Hiérarchie du layout :} Algorithme de groupement pour recréer la structure spatiale.
        \item \textbf{Génération de code :} Production de code \textbf{Flutter} multiplateforme .
    \end{enumerate}
\end{frame}

% --- SOLUTIONS APPORTÉES (Données) ---
\begin{frame}{Solutions apportées : Pipeline de données}
    \begin{itemize}
        \item \textbf{Nettoyage du dataset :} Création d'un pipeline de prétraitement pour corriger les erreurs du dataset Rico (éléments mal alignés, étiquettes incorrectes) .
        \item \textbf{Équilibrage :} Technique pour assurer une distribution uniforme des classes dans l'entraînement .
        \item \textbf{Résultats :} Précision moyenne (mAP) de \textbf{91,37\%} pour les éléments atomiques et \textbf{87,40\%} pour les conteneurs .
    \end{itemize}
\end{frame}

% --- LIMITES ---
\begin{frame}{Limites}
    \begin{itemize}
        \item \textbf{Taille des données :} Les jeux de données d'entraînement restent relativement petits malgré les bons résultats .
        \item \textbf{Cas complexes :} Difficultés persistantes pour certains agencements hiérarchiques spécifiques (ex: alignement complexe de cartes verticales).
        \item \textbf{Statisme :} Incapacité à distinguer certains types dynamiques comme la VIDEO de l'IMAGE sur des captures statiques.
        \item \textbf{Dépendance à la qualité :} Nécessite une vérification manuelle pour assurer un étiquetage complet lors de l'entraînement .
    \end{itemize}
\end{frame}

% --- CONCLUSION ---
\begin{frame}{Conclusion}
    \begin{itemize}
        \item Le modèle démontre que le \textbf{Deep Learning} est parfaitement adapté à la reconnaissance visuelle des composants GUI .
        \item L'approche à \textbf{deux modèles (YOLOv7)} surpasse les approches à modèle unique en raison de la diversité des échelles d'objets .
        \item Solution prometteuse pour réduire drastiquement le temps de développement front-end .
        \item \textbf{Futur :} Augmenter la diversité des données et affiner l'algorithme de groupement spatial .
    \end{itemize}
\end{frame}

\end{document}